\documentclass[10pt]{article}
\usepackage[margin=0.8in]{geometry}
\usepackage{amsmath}
\usepackage[colorlinks=true,linkcolor=blue,urlcolor=blue]{hyperref}
\setlength{\parindent}{0pt}
\setlength{\parskip}{3pt}
\title{Adaptime: Method Overview}
\author{Thesis project}
\date{Current protocol}

\begin{document}
\maketitle

\begin{abstract}
Adaptime studies whether small fitted adaptors can combine a time-series
foundation model with calendar-aligned retrieval while preserving the official
TIME test protocol. This note defines the fixed validation-selected family and
the causal rolling per-variate horizon Ridge; it reports no empirical result.
\end{abstract}

\section{Chronological protocol}

Each dataset and forecast term is divided into a fixed datastore, adaptation
training, adaptation validation, and the untouched official TIME test interval.
The first three regions end strictly before the test. Ridge candidates are fit
once on adaptation training, $K$ and $\alpha$ are selected on adaptation
validation, and the selected coefficients are frozen for test. There is no
train-plus-validation refit.

Let $H$ be the forecast horizon, $T_{\mathrm{test}}$ TIME's configured test
length, and $n=\lfloor T_{\mathrm{test}}/H\rfloor$ its official complete
test-window count per variate. Adaptation training contains $2n$ window
origins and validation contains $n$, giving three times the test-window count
before test. Their interval lengths are $H+(2n-1)s_f$ and
$H+(n-1)s_f$, where the prime stride $s_f$ depends on frequency rather than
on $H$: 127 for intraday/hourly, 27 for business-daily/daily, 11 for weekly,
5 for monthly, and 3 for quarterly data. The remaining earlier observations
form the datastore.

The ridge history length is the same as in the inherited vanilla TIME run for
its backbone: 8192 for Chronos-2, 4096 for TS-ICL, and 2048 for Chronos-Bolt.
TS-RAG retains its separate native length of 512 when used as an external
control. Adaptation training and validation retain only rows with this complete
fixed-length history, so their retrieval distances always compare equal-width
representations. The vanilla test pass instead uses every official window with
all history available at that date, capped at the backbone limit.

For a query time $t$, let $d_{\max}$ be the fixed datastore endpoint and $P$
the dataset period. The latest eligible neighbor date is
\begin{equation}
  d_t = \min(t-H,d_{\max})
        - \bigl(\min(t-H,d_{\max})-t\bigr) \bmod P.
\end{equation}
Earlier candidates walk backward from $d_t$ by a datastore stride that is a
multiple of $P$. This keeps the datastore fixed while aligning its phase to the
query.

By default every eligible earlier datastore date is accessible. An optional
global window cap is divided evenly across variates, rounded down, and retains
the most recent dates. The fitting set has an analogous optional recent-window
cap; global fitting and same-variate fitting are explicit alternatives. TS-RAG
requires every variate to retain at least 11
datastore dates. Shared preparation records these counts; before TS-RAG
representation extraction, its project-owned data adapter checks the minimum
globally, including when a previously completed artifact is reused. Generic
preparation remains available to the Ridge family, whose insufficient-history
path is the vanilla fallback described below.
An early training or validation origin that cannot supply the fixed history is
excluded from fitting, but no official test origin is removed. Virtual
$K=0$ is always admissible and needs no fit. Each positive-$K$ candidate must
separately exceed the configured training and validation support thresholds;
when no fitted candidate qualifies, TIME's unchanged test windows use cached
vanilla forecasts.

Let $q_t$ be the finite fraction of a query history. The default gate requires
$q_t\geq0.8$ before retrieval. Datastore histories satisfy the same threshold,
and their complete retrieved future must be finite. A query is RAG-eligible
only when it has at least $K_{\max}$ valid chronological neighbors, where
$K_{\max}=15$ for the default grid. Every candidate $K$ uses this common query
support and consumes the first $K$ neighbors from the same ordered list; fewer
than $K_{\max}$ valid neighbors makes the query ineligible for all candidates.
After selection is frozen, test extraction retrieves only the selected $K$;
a test query needs those $K$ neighbors rather than $K_{\max}$.

\section{Fixed ridge family}

For each query, $V$ is the vanilla foundation forecast and $C$ is the forecast
obtained when the $K$ retrieved neighbor trajectories are supplied as
covariates. For neighbor $i$, $Y_i$ is its observed horizon and $N_i$ is the
foundation model's vanilla forecast from that neighbor's own history. The full
design is
\begin{equation}
  X = [V,C,Y_1,\ldots,Y_K,N_1,\ldots,N_K].
\end{equation}
The paired $Y_i$ and $N_i$ terms expose neighbor residual information. Let
$s_{t,c}$ be the root mean squared seasonal difference over the complete
history before query $t$ for channel $c$. Only pairs whose two original dates
are finite contribute, and their count is the denominator. For
\texttt{full\_ridge\_shared}, one no-intercept ridge vector is shared across
variates and forecast horizons and fitted by MSSE:
\begin{align}
  \hat\beta_\alpha
    &= \arg\min_\beta \sum_{t,c,h}
       \left(\frac{(Y-V)_{t,c,h}-(X\beta)_{t,c,h}}{s_{t,c}}\right)^2
       + \alpha\|\beta\|_2^2, \\
  V' &= V + X\hat\beta_\alpha.
\end{align}

The primary setting is $K=10$, $\alpha=10^{-2}$. Validation searches
$K\in\{1,5,10,15\}$ and
$\alpha\in\{10^{-3},10^{-2},10^{-1}\}$ and compares every fitted method to
virtual $K=0$. The nested \texttt{cov\_ridge\_shared} design uses $[V,C]$;
\texttt{y\_ridge\_shared} uses $[V,Y_1,\ldots,Y_K]$; and
\texttt{full\_ridge\_per\_variate} retains the full design while fitting a
separate coefficient vector from each variate's training windows. All designs
retain the outer vanilla term $V$ in $V+X\beta$. Delta, convex, and native
multivariate variants are outside the current protocol.

Ridge fitting uses complete adaptation-training rows only; it never applies an
element-wise NaN mean to the normal equations. During validation and test, an
ineligible query receives the vanilla output for both $C$ and $V'$. This gate
uses only query and datastore information available at forecast time. Missing
future labels affect metric support only and cannot select the prediction path.

\section{Bayesian covariate baseline}

For each candidate $K$, every eligible adaptation-training window is one paired
trial. Let $z_i=1$ when the per-window MSSE of $C_i$ is lower than that of
$V_i$, $z_i=0$ when it is higher, and $z_i=1/2$ for a tie. With a
$\operatorname{Beta}(1,1)$ prior, the posterior-mean probability is
\begin{equation}
  p = \frac{1+\sum_i z_i}{2+n}.
\end{equation}
Adaptation validation compares these frozen candidate mixtures by MSSE and
selects their $K$, or the virtual vanilla candidate. The selected evidence is
then frozen before test. The soft baseline is
\begin{equation}
  B=(1-p)V+pC.
\end{equation}
An ineligible test window receives $V$ for both the hard covariate and Bayesian
branches. A task-level insufficient-training fallback fixes $p=0$ and makes
all family predictions equal to $V$.

A second Bayesian candidate calls Chronos-2 with every other variate supplied
as a past-only covariate. The same paired training-win,
$\operatorname{Beta}(1,1)$ posterior, validation selection, and vanilla
fallback logic mixes this forecast with $V$. No future covariate value is
exposed.

\section{Rolling per-variate horizon ridge}

The independent \texttt{rolling\_y\_ridge\_horizon} method has no validation
selection. At every official test query it fits one coefficient vector per
variate and horizon with design $[V,Y_1,\ldots,Y_K]$. Defaults are $K=15$ and
$\alpha=1$. Fitting dates come only from the query variate, are aligned to the
query's dataset-period phase, retain at most the latest 100 dates, and require
at least 64. Each fitting target must be fully observed before the test query.

Retrieval uses a causal cross-variate datastore at every fitting and query
date. Its default 10,000-window cap is divided evenly across variates, and the
smallest available balanced count fixes one common datastore size over all
adapted dates. Incremental sufficient statistics remove expired fitting dates
and add new ones as the test origin advances. A query with insufficient fitting
dates or retrieval support returns $V$.

\section{Evaluation}

The official-test evaluations cover vanilla $V$, hard retrieval-context $C$,
both Bayesian gates, the nested shared Ridge designs, shared and per-variate
full Ridge, the validation-selected fixed family, rolling horizon Ridge, and
the independently implemented TS-RAG control. Seasonal Naive supplies the task
scaling baseline. Each method receives its own standard TIME artifact on
identical official windows. The evaluator retains
MSE, MAE, RMSE, MAPE, sMAPE, MASE, ND, and CRPS. Seasonal Naive remains the
inherited external baseline used by broader TIME summaries. Missing dates are
never removed or shifted: forecast errors skip non-finite target dates and
MASE scales count only finite lagged pairs at their original timestamps. The
comparison reports each metric mean together with its finite and total value
counts; unequal finite coverage remains visible rather than preventing report
construction.
Conclusions require completed and inspected cluster artifacts.

Eligibility arrays, reason codes, effective training counts, and test RAG
coverage, Ridge coefficients, Bayesian evidence, rolling fitting support, and
cache identities are retained with the task artifacts. Vanilla forecasts and
exact-context representations are shared across fixed and rolling methods only
when prepared data, model, weights, source window, and context length agree;
neighbor selections remain method-specific. The resulting adapted estimands are transparent gated hybrids:
their selected method on eligible windows and vanilla otherwise. An
insufficient-training task is the limiting case with zero eligible Adaptime
windows and vanilla predictions for every fixed-family branch.

\end{document}
